本节只使用已登记的确定性几何计算和固定种子合成报告。所有合成案例均由同一输入脚本生成，策略只能读取与官方接口同结构的响应，不能访问源的真实位置、接收半径或发射方向。其作用是比较参数和算法版本，不替代正式测试成绩。

\subsection{第二检测点的参数敏感性}

式\eqref{eq:second-point}表明，第二检测点的横向基线由方位误差半宽
$\varepsilon$ 和距离比 $R/r_0$ 共同决定。固定 $r_0=1000\unit{m}$，分别改变一个参数，并对首次扇形中的可能源位置重新计算最坏交会角，结果见图\ref{fig:q2-sensitivity}。图中虚线对应题面参数
$\varepsilon=1^\circ$、$R/r_0=1.5$。

\begin{figure}[htbp]
  \centering
  \includegraphics[width=0.98\textwidth]{q2-parameter-sensitivity.pdf}
  \caption{第二检测点横向基线和最坏交会角的参数敏感性}
  \label{fig:q2-sensitivity}
\end{figure}

在固定 $R=1500\unit{m}$ 时，$\varepsilon$ 从 $0.25^\circ$ 增至
$3^\circ$，横向基线从 $658.159\unit{m}$ 降至 $621.279\unit{m}$，最坏交会角从 $41.098^\circ$ 降至 $35.576^\circ$。题面工作点的两个指标分别为 $648.248\unit{m}$ 和
$39.592^\circ$。该区间内变化连续，说明对 $1^\circ$ 附近的小幅误差扰动不敏感。

距离上界的影响更显著。$R/r_0$ 从 $1.0$ 增至 $1.95$ 时，横向基线从 $857.167\unit{m}$ 降至 $205.155\unit{m}$，最坏交会角从 $58.244^\circ$ 降至 $10.893^\circ$。当比例逼近 2 时，三圆盘公共区域沿横向收缩，交会几何迅速恶化。表\ref{tab:q2-sensitivity}列出关键节点。每一行推荐点到首次扇形中任意候选源的最大距离均不超过 $1000\unit{m}$；变化的是定位条件，而不是接收保证。

\begin{table}[htbp]
  \centering
  \caption{问题二关键参数点的稳健性指标}
  \label{tab:q2-sensitivity}
  \begin{tabular}{cccc}
    \toprule
    条件 & 横向基线/m & 最坏交会角/$(^\circ)$ & 最大候选源距离/m \\
    \midrule
    $\varepsilon=0.25^\circ,\ R/r_0=1.5$ & 658.159 & 41.098 & 1000.000 \\
    $\varepsilon=1.00^\circ,\ R/r_0=1.5$ & 648.248 & 39.592 & 1000.000 \\
    $\varepsilon=3.00^\circ,\ R/r_0=1.5$ & 621.279 & 35.576 & 1000.000 \\
    $\varepsilon=1.00^\circ,\ R/r_0=1.8$ & 420.116 & 23.958 & 1000.000 \\
    $\varepsilon=1.00^\circ,\ R/r_0=1.95$ & 205.155 & 10.893 & 1000.000 \\
    \bottomrule
  \end{tabular}
\end{table}

\subsection{覆盖路线与停止规则消融}

为分离两个策略改动的作用，使用随机种子 $20261001$ 对问题三、四各生成 50 个案例，并在完全相同的源布局上依次运行 v1、v2、v3。v1 在定位绕行后继续固定蛇形顺序；v2 改为前往最近的未完成覆盖点；v3 在 v2 基础上加入“已清除 16 个源立即退出”。三版保持同一 61 点覆盖集、示向区域更新和清除判据。

\begin{table}[htbp]
  \centering
  \caption{三个策略版本的同种子合成消融结果（每问50局）}
  \label{tab:ablation}
  \resizebox{0.92\textwidth}{!}{%
  \begin{tabular}{ccrrrr}
    \toprule
    问题 & 版本 & 全部清除局数 & 平均虚拟时间/s & 标准差/s & 平均移动距离/m \\
    \midrule
    3 & v1 & 50/50 & 18148.492 & 620.599 & 69226.362 \\
    3 & v2 & 50/50 & 15843.685 & 585.260 & 62092.623 \\
    3 & v3 & 50/50 & 15430.448 & 1574.642 & 60324.542 \\
    \midrule
    4 & v1 & 50/50 & 18158.743 & 492.772 & 68224.216 \\
    4 & v2 & 50/50 & 16344.127 & 642.935 & 62386.034 \\
    4 & v3 & 50/50 & 16092.024 & 1243.202 & 61308.018 \\
    \bottomrule
  \end{tabular}}
\end{table}

由 v1 到 v2，问题三、四平均虚拟时间分别下降 $12.70\%$、
$9.99\%$，平均移动距离分别下降 $10.30\%$、$8.56\%$，说明主要收益来自吸收定位绕行后的返程。v3 的提前停止只会在源数恰为题面上界 16 时触发；50 局中每问各有 4 局属于该子集。在这 4 组配对案例中，问题三平均时间由 $15848.940\unit{s}$ 降至 $10683.487\unit{s}$，下降 $32.59\%$；问题四由 $16078.081\unit{s}$ 降至 $12926.794\unit{s}$，下降 $19.60\%$。图\ref{fig:ablation}把全样本路线收益与 16 源停止收益分开展示。

\begin{figure}[htbp]
  \centering
  \includegraphics[width=0.98\textwidth]{q34-strategy-ablation.pdf}
  \caption{同种子合成案例中的路线与停止规则消融}
  \label{fig:ablation}
\end{figure}

v3 在全样本上的标准差高于 v2，是因为 16 源案例被大幅提前截断，形成与 10--15 源案例不同的时间层，而不是运行稳定性下降。所有展示组的清除比例仍为 1。由于 v3 没有改变 10--15 源分支，不能把全样本均值的下降外推为任意源数上的普遍加速。

\subsection{定向源比例压力分析}

问题四的 50 局合成样本含 10--16 个源，定向源比例覆盖
$33.33\%$--$90.91\%$。按比例小于 $50\%$、$50\%$--$70\%$、不小于 $70\%$ 分层后，样本量分别为 11、26、13，平均虚拟时间分别为 $15256.308\unit{s}$、$16197.544\unit{s}$、$16588.126\unit{s}$，三组均全部清除。图\ref{fig:directional-stress}左侧同时用颜色标出总源数，避免把源数差异遗漏在比例比较之外。

\begin{figure}[htbp]
  \centering
  \includegraphics[width=0.98\textwidth]{q4-directional-stress.pdf}
  \caption{问题四合成案例的定向源比例压力分析}
  \label{fig:directional-stress}
\end{figure}

分层均值随定向比例上升，但本实验没有控制总源数、具体位置和发射方向，且各层样本量不同，因此只能说明测试集覆盖了高定向比例情形，不能据此认定定向比例对耗时存在单独的因果效应。可以确认的是，在该压力范围内只使用正示向更新、保留无信号候选位置的保守策略没有出现漏清除。

\subsection{分析证据的可复现性}

本节所有数值由 \texttt{python/b\_paper\_analysis.py} 从三份已登记报告派生，结果固化在
\texttt{output/paper-evidence-20260912.json}。生成脚本不重新随机抽样，并逐问题、逐种子核对三个版本的源数和完整合成真值，输入不一致时直接报错；相同输入会写出字节稳定的 JSON。自动测试同时检查配对一致性、题面工作点、两个关键降幅、50 局分层计数和全部清除状态。图表脚本只读取该 JSON，不在绘图阶段重新计算或修改实验结论。
